\section{描述性统计分析}

基于本次杭甬温三市环境噪声感知调研的467份有效问卷数据，本次调研从受访者样本特征、噪声暴露与感知整体情况、声环境认知现状等维度开展全面的描述性统计分析，系统呈现调研样本的结构特征、居民噪声暴露与感知的整体态势，以及对声环境治理的认知与需求情况。

\subsection{受访者样本特征分析}

\subsubsection{人口统计学特征}

\begin{figure}[H]
\centering
\begin{subfigure}[b]{0.46\textwidth}
  \centering
  \includegraphics[width=\textwidth]{板块1_性别构成.png}
  \caption{性别构成}
\end{subfigure}
\hfill
\begin{subfigure}[b]{0.46\textwidth}
  \centering
  \includegraphics[width=\textwidth]{板块2_城市分布.png}
  \caption{城市分布}
\end{subfigure}

\vspace{0.6em}

\begin{subfigure}[b]{0.46\textwidth}
  \centering
  \includegraphics[width=\textwidth]{板块3_年龄构成.png}
  \caption{年龄构成}
\end{subfigure}
\hfill
\begin{subfigure}[b]{0.46\textwidth}
  \centering
  \includegraphics[width=\textwidth]{板块4_文化程度.png}
  \caption{文化程度}
\end{subfigure}

\vspace{0.6em}

\begin{subfigure}[b]{0.46\textwidth}
  \centering
  \includegraphics[width=\textwidth]{板块5_居住区域.png}
  \caption{居住区域}
\end{subfigure}
\hfill
\begin{subfigure}[b]{0.46\textwidth}
  \centering
  \includegraphics[width=\textwidth]{板块6_职业构成.png}
  \caption{职业构成}
\end{subfigure}
\caption{受访者样本人口统计学特征\label{fig:profile}}
\end{figure}

\textbf{（1）性别与年龄特征}

如图\ref{fig:profile}(a)、(c)所示，样本中男性260人（55.67\%），女性207人（44.33\%），男性比例略高于女性。年龄结构以中青年为主体，31-45岁群体比例最高（31.69\%），19-30岁（27.41\%）及46-60岁（22.70\%）次之，18岁以下（8.35\%）及61岁及以上（9.85\%）比例相对较低。

\textbf{（2）城市与居住区域分布}

如图\ref{fig:profile}(b)、(e)所示，样本数据采集自杭州、宁波、温州三市，分布与前期调研权重设计一致。其中宁波市样本量最大（192人，41.11\%），温州市（143人，30.62\%）与杭州市（132人，28.27\%）次之。居住区域类型方面，城市主城区（134人，28.69\%）与城市郊区（132人，28.27\%）占比相近，县城（90人，19.27\%）、乡镇（74人，15.85\%）与农村（37人，7.92\%）依次递减，样本以城区居民为主，同时保留了相当比例的乡镇与农村样本。

\textbf{（3）受教育程度}

如图\ref{fig:profile}(d)所示，样本中本科及以上学历者合计超过四成，其中本科学历160人（34.26\%），硕士及以上31人（6.64\%）。此外，高中/中专/技校（19.27\%）、大专（17.34\%）、初中（14.56\%）、小学及以下（7.92\%）等不同学历层次均有覆盖。

\textbf{（4）职业分布}

如图\ref{fig:profile}(f)所示，样本职业类型较为多元。企业工作人员比例最高（130人，27.84\%），学生群体次之（119人，25.48\%）。自由职业者/个体工商户（65人，13.92\%）、农林牧渔从业者（38人，8.14\%）等亦占有一定比例，基本覆盖了城市主要就业群体。

\subsubsection{生活状态特征}

\fig{图片11.png}{隔音墙状态}

居住隔音配套是决定居民噪声暴露水平的关键客观硬件条件，本次调研专门设置住宅隔音墙相关题项统计现状。数据显示仅28\%受访者家中配备专业隔音墙体或隔音门窗配套，高达72\%的住宅无隔音设施或受访者无法确认房屋隔音配置。住宅隔音设施配套失衡是老旧小区改造遗留的共性短板，新建商品房隔音门窗普及率更高，而建成10年以上老旧楼栋几乎无专业隔音构造，这一硬件差距造成居民日常噪声暴露强度分层，也与老旧小区占比偏高的特征形成互证。老旧小区墙体单薄、窗户密封性差，导致外部交通、施工噪声可直接穿透进入室内，大幅加剧居民主观困扰程度，也是后续治理需要重点补齐的硬件短板。

\subsection{噪声暴露与感知整体情况}

\fig{图片12.png}{噪声暴露全景分析}

\subsubsection{噪声困扰程度分布}

本研究通过李克特5级量表对噪声影响程度进行评估。

结果显示：认为噪声“影响一般”与“影响较大”的受访者占比并列最高，均为30.62\%；其次为“偶尔有影响”（22.70\%）、“几乎没有影响”（11.35\%）和“影响非常大”（4.71\%）。超60\%的受访者认为噪声对生活存在中等及以上影响，其中超三成居民受到较大程度干扰，反映出杭甬温三市居民的噪声困扰问题较为普遍，噪声污染已成为影响居民生活质量的重要环境问题。

\subsubsection{各类噪声源干扰特征}

从干扰程度排序来看，交通噪声（3.09分）与施工噪声（3.01分）的干扰强度最高，二者得分均超过3分，处于中等偏上水平；其后依次为动物噪声（2.95分）、公共活动噪声（2.83分）、生活噪声（2.76分）与商业噪声（2.70分）。整体而言，交通噪声与施工噪声是杭甬温三市居民面临的最主要噪声干扰源，这与城市交通流量大、工程建设活动频繁的发展特征高度契合。

\subsubsection{噪声影响时段与距离特征}

从时段特征来看，噪声影响呈现明显的夜间集中倾向，晚上和深夜受到噪声干扰的受访者数量最多，分别为138人和135人，白昼人数占比较少。这一分布特征表明，夜间噪声对居民的影响更为突出，与居民休息、睡眠的核心时段高度重合，更容易对生活质量造成负面影响。

从声源距离分布来看，认为噪声源距离适中的人数最多，对应161人；其次为距离较近（102人）和较远（101人）；近在咫尺（36人）和非常远（67人）的受访者占比相对较低。这一结果说明，大部分居民面临的噪声源来自周边50米范围内的环境，以社区及周边活动、道路交通等近距离声源为主，这类噪声更易对居民生活产生持续且直接的干扰。

整体来看，噪声影响的时空特征与居民的生活作息和居住环境高度相关，夜间与近距离噪声是当前治理工作中需要重点关注的方向。

\subsection{声环境认知现状}

\fig{图片13.png}{满意度总览}

\subsubsection{声环境与治理满意度评价}

对于声环境质量满意度，超44\%的受访者对当前的声环境质量持不满意态度，仅约27\%的受访者持满意态度，反映出杭甬温三市居民对当前声环境质量的整体评价偏低，声环境质量与居民的预期仍存在较大差距。

对于噪声治理工作满意度，约40\%的受访者对当前的噪声治理工作持满意态度，超30\%的受访者持一般态度，仅约30\%的受访者持不满意态度，居民对噪声治理工作的满意度显著高于对声环境质量的满意度，反映出居民对本地噪声治理工作的努力有一定的认可，但治理成效仍未达到居民的预期，有较大的提升空间。现有管控力度、覆盖范围、整改处置效率尚未完全消解长期噪声侵扰带来的负面体验，治理成效存在明显滞后性，管控行动转化为实际环境仍存在堵点。

\subsubsection{居民应对方式分析}

\fig{图片14.png}{面对噪声的应对方式}

整体来看，居民遭遇噪声污染时，最倾向于采用自主协商、向社区/物业反馈的私力救济与基层救济方式，同时自我防护也是重要的应对手段。而官方渠道投诉、法律途径维权等公力救济方式的采用比例相对较低，仅约40\%的受访者会选择官方投诉渠道，不足20\%的受访者会采取法律途径维权，反映出当前居民的噪声维权渠道的使用意愿仍有较大的提升空间，维权渠道的便捷性、有效性仍需进一步优化。

\subsubsection{噪声治理认知与需求}

\fig{图片15.png}{噪声治理倾向\label{fig:tendency}}

由图\ref{fig:tendency}可知，超61\%的受访者认为强化浙江省环境噪声治理有必要或非常有必要，仅约18\%的受访者认为没有必要或完全没有必要。这说明绝大多数居民对噪声治理的必要性有高度的认知，具备较强的噪声治理意愿与需求，为噪声治理工作的推进提供了良好的群众基础。

\fig{图片16.png}{噪声治理需求\label{fig:demand}}

由图\ref{fig:demand}可知，居民对噪声治理的需求呈现出“制度完善优先、源头管控核心、硬件保障配套、监管执法兜底”的特征，最核心的需求是完善噪声管控的法律法规与标准，体现现有管控细则仍存在模糊地带，基层执法缺少清晰量化分贝、时限依据。

其次居民对施工、交通噪声的时段管控、隔音硬件设施建设、投诉处理机制优化等方面也有强烈的需求，群众更期待兼顾制度完善与实体降噪工程同步推进，这为后续噪声治理工作的方向与重点提供了清晰的民意依据。

